\subsection{Scaffold semantics mediated real-data amplification}

Changing only the acyclic-scaffold meaning strongly attenuated the empirical regression pattern. Under singleton semantics, ESOL mean effects declined to 0.036--0.061 RMSE units, while all three FreeSolv point estimates reversed sign and all intervals crossed zero (Figure~\ref{fig:acyclic-primary}). The empirical--null bridge independently showed that singleton FreeSolv no longer departed from random endpoint assignment. Together these findings identify acyclic-scaffold semantics as a benchmark hyperparameter rather than a cosmetic implementation detail.

\subsection{Molecular-record policy altered directional conclusions}

The source-faithful data contained 105 multi-component BBBP records, 14 ClinTox records, and 3,086 HIV records. Deterministic dominant-fragment selection changed scaffold identity for 5, 1, and 235 records, respectively, and created new duplicate and conflicting-label identities. All nine corrected classification sensitivity inferences remained inconclusive, but seven of nine mean-effect directions reversed (Figure~\ref{fig:representation}). Thus, inferential status was more stable than small directional point estimates, while molecular-record policy remained part of the benchmark specification.

\begin{figure}[!htbp]
\centering
\includegraphics[width=0.96\textwidth]{figures/Figure_6.pdf}
\caption{Molecular-record representation sensitivity. Panel A summarizes structural consequences of deterministic dominant-fragment mapping. Panel B connects source-faithful and dominant-fragment mean AUC effects for the nine affected dataset--model cells. Corrected inferences remained inconclusive in both representations, but seven mean-effect directions reversed.}
\label{fig:representation}
\end{figure}

\subsection{Protocol-conditioned empirical effects accompanied the mechanism}

Response-aware selection changed benchmark properties beyond the optimized endpoint mean. The mean-only effect exceeded the corresponding learned-model mean effect in all six primary regression cells. Across datasets, mean paired changes ranged from $-0.060$ to 0.015 in largest-test-scaffold fraction and from $-10.993$ to 48.752 in effective test scaffold number. The optimized gap necessarily decreased, but the selected molecular population also moved in structural and response-distribution dimensions that were not fixed by the paired design. Figure~\ref{fig:collateral} places the mean-only control beside these collateral changes.

\begin{figure}[!htbp]
\centering
\includegraphics[width=0.96\textwidth]{figures/Figure_7.pdf}
\caption{Protocol-conditioned empirical effects and collateral diagnostics. Panel A compares the mean-only regression control with learned-model effects. Panel B shows the paired reduction in target-mean mismatch. Panels C and D show changes in test-set scaffold concentration and effective scaffold number, neither of which entered the response-aware selection objective.}
\label{fig:collateral}
\end{figure}
